\section{Fazit}

Die Arbeit hat untersucht, wie das Modell Pythia-410M Sentiment intern verarbeitet. Als Fallstudie der \gls{mechanisticinterpretability} wurden alle fünf Werkzeuge angewendet, die von beschreibenden hin zu kausalen Aussagen führen. Die Untersuchung folgte einem durchgehenden roten Faden von der Frage, ob Sentiment vorhanden ist, über die Frage, ab wann das Modell es nutzt, bis zur Frage, welche Komponenten beteiligt sind.

Die zentralen Befunde lassen sich knapp zusammenfassen. Sentiment ist bereits im Embedding-Raum strukturiert und linear zu 97 Prozent trennbar. Über die Schichten hinweg wird die Trennung zwischen positiven und negativen Eingaben schrittweise aufgebaut und erreicht ihr Maximum in den oberen Schichten zwischen etwa Schicht 17 und Schicht 23. Mehrere \glspl{attentionhead} richten ihre Aufmerksamkeit auf sentimenttragende Wörter, besonders in den frühen bis mittleren Schichten. Das Activation Patching zeigt für ein einzelnes Prompt-Paar, dass ein Eingriff am Stimmungswort die Vorhersage tatsächlich wiederherstellt, allerdings an einer früheren, token-nahen Stelle als von Logit Lens und \gls{attention}-Analyse nahegelegt und ohne dass sich ein einzelner verantwortlicher Kopf identifizieren lässt. Auf der Verhaltensebene hängt die Sentiment-Erkennung stark vom Prompt-Design ab, und komplexe Phänomene wie Sarkasmus bleiben schwierig. Diese Befunde stützen den beschreibenden und für ein Beispiel auch den kausalen Teil von Hypothese H1 sowie den beobachtbaren Teil von Hypothese H3. Hypothese H2 bestätigt sich nur teilweise: Die Position ist scharf lokalisiert, die Schicht nicht.

Der nächste und entscheidende Schritt ist, das Activation Patching über mehrere oder alle 1.707 Paare des CAD-Datensatzes zu aggregieren, so wie es für Logit Lens und \gls{attention}-Analyse bereits geschehen ist. Erst eine solche Aggregation kann zeigen, ob die frühe, token-nahe Lokalisierung aus dem Einzelbeispiel verallgemeinerbar ist oder ob sie von der konkreten Formulierung abhängt, und schließt damit die Argumentationskette von der Korrelation bis zur Kausalität endgültig. Darüber hinaus bieten sich eine Erweiterung auf größere Modelle der Pythia-Reihe, der Einsatz reichhaltigerer Sentiment-Maße jenseits eines festen Lexikons sowie eine genauere Untersuchung der schwierigen sprachlichen Kategorien als weiterführende Arbeiten an.
